\chapter{مقدمه}
\label{ch:intro}

این نوشتار \مهم{نسخه خلاصه} گزارش پروژه کارشناسی است. گزارش کامل، با هفت
فصل، واژه‌نامه و پیوست بازتولید، در پرونده جداگانه‌ای آمده است. اعداد،
نیازمندی‌ها و مرز ادعاها در هر دو نسخه یکی است؛ آنچه کوتاه شده، شرح و
تکرار است نه نتیجه.

گفت‌وگوی گفتاری طبیعی‌ترین شکل تعامل انسان با ماشین است، اما تا همین اواخر
تقریباً همیشه از زنجیره‌ای از مؤلفه‌های مستقل ساخته می‌شد. مدل‌هایی مانند
\en{Moshi}~\cite{defossez2024moshi} و
\en{Freeze-Omni}~\cite{wang2025freezeomni} نشان دادند که می‌توان گفت‌وگو را
به‌صورت یک مسئله «گفتار‌به‌گفتار» و در قالب مدلی که هم‌زمان می‌شنود و سخن
می‌گوید مدل کرد. این تحول برای زبان‌های پرمنبع رخ داده است. برای فارسی،
هیچ‌یک از مدل‌های متن‌باز موجود پشتیبانی رسمی ارائه نمی‌کنند و پیکره
مکالمه‌ای با برچسب وقفه نیز به‌صورت عمومی در دسترس نیست.

پرسش عملی پروژه این است: با بودجه یک پروژه کارشناسی، تا چه اندازه می‌توان
یک سامانه گفتار‌به‌گفتار فارسی تمام‌دوطرفه ساخت، و چه چیزی را می‌توان با
شواهد قابل ممیزی «اثبات» کرد؟

\section{تعریف مسئله}
\label{sec:problem}

هدف تعریف‌شده، طراحی، پیاده‌سازی و ارزیابی یک نمونه اولیه مدل زبانی بزرگ
گفتار‌به‌گفتار برای زبان فارسی است. سامانه باید تمام‌دوطرفه باشد: میکروفون
کاربر در طول پخش پاسخ فعال بماند و سامانه بتواند به «وقفه» کاربر واکنش
نشان دهد. سه ویژگی این مسئله را از پرسش‌و‌پاسخ گفتاری متعارف جدا می‌کند:
هم‌زمانی شنیدن و گفتن؛ تصمیم توقف پخش در مقیاس چند‌ده میلی‌ثانیه؛ و
تداوم گفتاری که در لحظه وقفه ادا می‌شود، تا کاربر ناچار به تکرار جمله
نشود.

\section{اهمیت موضوع}

رابط گفتاری برای کاربرانی که با صفحه‌کلید فارسی راحت نیستند، برای
کم‌بینایان و به‌طور خاص برای سالمندان می‌تواند تنها راه مؤثر استفاده از
خدمات دیجیتال باشد. برای فارسی، منابع بازشناسی و تبدیل متن به گفتار وجود
دارد؛ نمونه شاخص \en{ManaTTS}~\cite{qharabagh2025manatts} با بیش از
\N{86} ساعت گفتار تک‌گوینده است. اما این منابع خوانشی‌اند و هم‌پوشانی،
وقفه و بازخورد کلامی را پوشش نمی‌دهند. از سوی دیگر، مدل‌های روز مانند
\en{Qwen2.5-Omni}~\cite{xu2025qwen25omni} برای نسخه هفت‌میلیارد‌پارامتری
خود به حافظه‌ای به‌مراتب بیش از بازه \N{12} تا \N{24} گیگابایت نیاز دارند.
نشان‌دادن آنچه در این بازه قابل دستیابی است و آنچه نیست، خودش یک نتیجه
مهندسی است.

\section{نیازمندی‌های تعریف‌شده}
\label{sec:requirements}

تعریف پروژه هفت نیازمندی مشخص را تعیین کرده است
(جدول~\ref{tab:requirements}). خروجی مورد انتظار نیز چهار قلم است: نمونه
اولیه کارا، کد منبع مستند، مجموعه داده مکالمه فارسی، و یافته‌های مربوط به
تعامل سالمندان. برای قلم چهارم، پروتکل و زیرساخت سنجش آماده شده است؛ یافته
تجربی آن تنها پس از تکمیل مطالعه انسانی قابل ادعا خواهد بود و ازاین‌رو در
این نوشتار ادعایی درباره تجربه سالمندان مطرح نمی‌شود.

\begin{table}[ht]
\centering
\caption{نیازمندی‌های هفت‌گانه تعریف پروژه}
\label{tab:requirements}
\small
\begin{tabular}{clL{6.4cm}}
\toprule
\textbf{شماره} & \textbf{نیازمندی} & \textbf{سنجه تعیین‌شده} \\
\midrule
ن۱ & نمونه اولیه گفتار‌به‌گفتار فارسی & تولید پاسخ گفتاری فارسی از ورودی
      گفتاری فارسی \\
ن۲ & کارکرد تمام‌دوطرفه & شنیدن در هنگام پخش و واکنش به وقفه \\
ن۳ & تأخیر سرتاسری & حداکثر \N{500} میلی‌ثانیه از دریافت گفتار تا تولید پاسخ
      گفتاری \\
ن۴ & مجموعه داده مکالمه فارسی & \N{100} تا \N{200} ساعت گفتار همراه با نویز،
      هم‌پوشانی و قطع‌و‌وصل کلام \\
ن۵ & ماژول تشخیص وقفه & دقت بالای \N{80} درصد با ویژگی‌های انرژی، زیر‌و‌بمی و
      \en{MFCC} \\
ن۶ & ارزیابی روی سخت‌افزار هدف & پردازنده گرافیکی با \N{12} تا \N{24} گیگابایت
      حافظه \\
ن۷ & ارزیابی انسانی & \N{5} تا \N{10} گویشور فارسی، دست‌کم دو نفر بالای
      \N{60} سال \\
\bottomrule
\end{tabular}
\end{table}

\section{راهبرد دومسیره و چالش‌ها}
\label{sec:two-track}

چهار چالش شکل نهایی کار را تعیین کرد: نبود مسیر آموزشی رسمی برای فارسی
در مدل‌هایی مانند \en{Mini-Omni2}~\cite{xie2024miniomni2} و
\en{LLaMA-Omni2}~\cite{fang2025llamaomni2}؛ نبود پیکره مکالمه‌ای
برچسب‌دار فارسی؛ فاصله میان کاهش زیان آفلاین و تولید خودرگرسیو معتبر؛ و
مرز میان شواهد مهندسی و ادعای علمی.

برای کاهش ریسک، دو مسیر به‌طور موازی پیش رفت. مسیر مستقیم، تطبیق
کم‌پارامتر \en{Moshika} با آموزش‌دهنده رسمی خانواده
\en{Moshi}~\cite{moshifinetune2024} است. مسیر آبشاری، زنجیره بازشناسی
گفتار فارسی، مدل زبانی متنی و تبدیل متن به گفتار است که کنترل تمام‌دوطرفه
و منطق وقفه به‌صورت جداگانه روی آن سوار می‌شود. هر دو مسیر از یک مجموعه
داده مشترک استفاده کردند. نتیجه نامتقارن بود: مسیر آبشاری نمونه اولیه
تحویل‌شده شد و مسیر مستقیم به‌عنوان نتیجه منفی مستند، با یک سیگنال
یادگیری مثبت، گزارش شد.

\section{دستاوردها}
\label{sec:contributions}

\begin{enumerate}
\item \مهم{نمونه اولیه کارکردی با کنترل تمام‌دوطرفه.} جریان پیوسته
      \en{PCM16} از مرورگر دریافت می‌شود، آشکارساز غلتان بر پایه انرژی،
      زیر‌و‌بمی و \en{MFCC} درباره وقفه تصمیم می‌گیرد، منبع پخش متوقف
      می‌شود و پیش‌بافر \N{450} میلی‌ثانیه‌ای به نوبت بعدی منتقل می‌گردد.
      در پنل ثابت اعتبارسنجی، زنجیره با پاسخ‌گوی \en{Qwen2.5-0.5B} در هر
      \N{9} ردیف رونویسی فارسی، پاسخ فارسی و صوت غیرساکت تولید کرد و به
      پاسخ جایگزین قاعده‌محور بازنگشت. پاسخ‌گوی تحویل‌شده سپس به
      \en{Qwen3-4B} با پرامپت \en{v2} ارتقا یافت؛ آن پنل پس از ارتقا دوباره
      اجرا نشد.

\item \مهم{مجموعه داده ممیزی‌شده.} از موجودی \N{775.887} ساعتی،
      \N{6,754} جفت با \N{108.584} ساعت صوت استریو ساخته شد که درون بازه
      \N{100} تا \N{200} ساعت است. تفکیک بر پایه شناسه نشست منبع انجام شد و
      نشت گروهی صفر بود.

\item \مهم{نتیجه مثبت خودکار معنایی.} پس از صلاحیت هفت‌ردیفی، آزمون نهایی
      \N{40} ردیفی از هر هفت شرط خودکار عبور کرد: ارتباط از \N{1.425} به
      \N{3.150} و انسجام به \N{3.325} رسید، با نرخ برد زوجی \N{67.5} درصد و
      اعتبار \N{240/240} فراخوانی داور. داور هم‌خانواده است و این سنجه
      جایگزین ارزیابی انسانی نیست.

\item \مهم{نتیجه منفی مستند برای مسیر مستقیم.} شش نسل آزمایشی زیان را کاهش
      داد اما هیچ چک‌پوینتی از شرط اجرای خودرگرسیو عبور نکرد. نسخه
      \en{v6.2} کاهش \N{32.20} درصدی زیان متنی درون‌نمونه داشت و هر چهار
      چک‌پوینت آن \N{0} از \N{9} گرفتند.

\item \مهم{روش کاری شواهد‌محور.} پروتکل‌ها پیش از اجرا قفل شدند، مصنوعات
      یک‌بار‌نوشت‌اند، و تجمیع‌کننده شکست‌بسته مانع از ارتقای تقریب خودکار
      به شواهد رسمی می‌شود.
\end{enumerate}

برای چهار نیازمندی هنوز شواهد رسمی کامل گردآوری نشده است: تأخیر سرتاسری
\N{500} میلی‌ثانیه، عبور رسمی تشخیص وقفه از \N{80} درصد با برچسب مستقل،
ارزیابی روی پردازنده گرافیکی فیزیکی هدف، و مطالعه انسانی.

\section{ساختار این نوشتار}

فصل~\ref{ch:background} مفاهیم لازم و جایگاه کار را فشرده می‌کند.
فصل~\ref{ch:method} روش ساخت داده، معماری آبشاری، آشکارساز وقفه، آزمایش‌های
مسیر مستقیم و کلاس‌های شواهد را شرح می‌دهد. فصل~\ref{ch:results} نتایج
اندازه‌گیری‌شده را با مرز ادعای هر یک گزارش می‌کند.
فصل~\ref{ch:conclusions} تفسیر، محدودیت‌ها و کارهای آینده را جمع می‌بندد.
